\appendix

\noindent This appendix provides complete technical foundations, detailed proofs, and practical guidance for the hierarchical summarization framework. It is organized to support both theoretical understanding and empirical deployment.

\medskip
\noindent\textbf{Appendix~\ref{app:notation}: Notation Recap and Local Foundations} collects all notation, definitions, and the foundational Lemma~\ref{lem:nodewise} used throughout. Readers comfortable with the main text may skip to Appendix~\ref{app:proofs-core}.

\noindent\textbf{Appendix~\ref{app:proofs-core}: Proofs of Main Results} establishes schedule invariance (Corollary~\ref{cor:schedule}), fold-of-folds invariance (Corollary~\ref{cor:folds}), and multi-round preservation (Theorem~\ref{thm:multi-round})—the three core guarantees that make hierarchical processing practical.

\noindent\textbf{Appendix~\ref{app:blackwell}: Blackwell-Style Score Transport} shows how supervision signals factor through the oracle $\fstar$ under Assumption~\ref{ass:CF}, enabling training on summaries instead of full documents.

\noindent\textbf{Appendix~\ref{app:practical_audit}: Practical Audit Guidance} provides step-by-step audit procedures for deployed hierarchies, including sample-size calculations and methods for composing leaf-level and merge-level distortions into deployment-level bounds.

\noindent\textbf{Appendix~\ref{app:bootstrap}: Local Consistency Bootstrap} describes the DSPy-style optimization loop that alternates between sampling local checks and updating the prompt parameters until the empirical violation rates fall below a desired threshold.

\noindent\textbf{Appendix~\ref{sec:dpo}: Direct Preference Optimization} connects oracle preservation to preference learning, proving that DPO on summaries recovers the same optimal policies as DPO on full documents when the sufficiency, merge-consistency, and idempotence conditions hold.

\noindent\textbf{Appendix~\ref{sec:global}: Global Compatibility Variant} presents an alternative formulation using global merge-consistency/associativity assumptions, connecting the framework to algebraic structures like monoid congruences and mergeable sketches.

\noindent\textbf{Appendix~\ref{sec:info-sieve}: Information-Theoretic View} (optional) sketches how OPS can be interpreted through an information-decomposition lens akin to the Information Sieve \citep{steeg2016informationsieve}, though this perspective is not required for the main arguments.

\noindent\textbf{Appendix~\ref{app:existence}: Existence of Oracle-Preserving Summaries} establishes that oracle-preserving summarizers always exist (set-theoretically) and can be implemented as compact encodings when the oracle space admits short representations.

\noindent\textbf{Appendix~\ref{sec:examples}: Worked Examples} instantiates the framework for binary event detection, entity extraction, counting aggregates, and provides a counterexample illustrating why stability cannot be omitted.

\medskip


\section{Notation Recap and Local Foundations}
\label{app:notation}

This appendix collects the primitives needed before the first theorem in Section~\ref{sec:setup}. It spells out every object once, keeping the notation flexible enough that we can later drop explicit ``$g$" symbols and simply write $g$ for the hierarchical evaluation.

\subsection{Strings, oracle, and metric}
We work on the free monoid $(\Strings,\concat,\emptystr)$: $\Strings$ is the set of all finite token sequences, $\concat$ is the usual concatenation of strings, and $\emptystr$ is the empty sequence acting as the identity. Documents are random variables $X\in\Strings$ drawn from a population law $\rho$. When we fix a specific document, we denote its base, unsummarized form by $x_0$ so that the realized leaf blocks are literally contiguous substrings of $x_0$. A fixed oracle $\fstar:\Strings\to Y$ maps strings to their task value (labels, tuples, rewards, etc.). A measurable (pseudo)metric $\metric:Y\times Y\to[0,\infty)$ scores task disagreement, and we use it directly as $\metric(\fstar(z),\fstar(x))$.

\subsection{Summarizers and expectations}
A summarizer is a (possibly stochastic) Markov kernel $g:\Strings\rightsquigarrow\Strings$ that turns any input string into a string-valued random output. Whenever randomness is present we average over the seeds or call counts used by $g$; when $g$ is deterministic the equalities below hold pointwise. To avoid notational clutter we simply write $\E[\cdot]$ for all expectations and suppress the implicit conditioning on the realized substrings.

\subsection{Partitions, trees, and hierarchical evaluation}
Fix a document $x=b_1\concat\cdots\concat b_k$; equivalently $x_0=x$ and $b_1,\ldots,b_k$ are contiguous substrings drawn from $x_0$. A \/emph{realized partition and tree} consists of those blocks together with a full binary tree $T$ whose leaves, read left to right, are $(b_1,\ldots,b_k)$. Leaves therefore contain no summaries—only raw text—and internal nodes contain summaries with two ordered children denoted $u_L,u_R$. Each node $u\in T$ carries two associated strings. The latent raw span is
\[
S(u)=
\begin{cases}
b_i,& u\text{ is the $i$th leaf},\\
S(u_L)\concat S(u_R),& u\text{ internal with children }u_L,u_R,
\end{cases}
\]
and the stored summary that flows upward is
\[
g(u):=
\begin{cases}
g\big(S(u)\big),& u\text{ leaf},\\
g\big(g(u_L)\concat g(u_R)\big),& u\text{ internal}.
\end{cases}
\]

\paragraph{Summary of symbols and expectations.}
The following table collects the core notation used throughout the appendix and main text.

\begin{table}[h]
\centering
\begin{tabular}{ll}
\toprule
Symbol & Meaning \\
\midrule
$x_0$ & Base (unsummarized) document; each leaf block $b_i$ is a contiguous substring \\
$b$ & A realized leaf block input to $g$ \\
$u$ & Internal node with children $u_L,u_R$; raw text $S(u)=S(u_L)\concat S(u_R)$ \\
$r$ & Root of tree $T$ \\
$g(T')$ & Summary at root of subtree $T'$; equivalently $g(T')=g(r')$ for root $r'$ of $T'$ \\
$Z^{(R)}(x)$ & $R$th-round summary: $Z^{(1)}(x)=g(r)$, $Z^{(R+1)}(x)=g(Z^{(R)}(x))$ \\
$\E$ & Expectation over randomness in $g$ (pointwise if $g$ deterministic) \\
\bottomrule
\end{tabular}
\caption{Core notation for hierarchical summarization.}
\label{tab:notation}
\end{table}

\subsection{Consistency conditions}
The guarantees in the main text follow from three statements checked \emph{only} on the leaves and internal nodes that the realized tree touches. All three definitions coincide with their counterparts in Section~\ref{sec:setup}; we restate them here to make the appendix self-contained.

\begin{definition}[Sufficiency condition (C1)]
For each realized leaf $b$, the single-block summary preserves the oracle in expectation:
\[
\E\big[\metric\big(\fstar(g(b)),\fstar(b)\big)\big]=0.
\]
This is Eq.~\eqref{eq:leaf_sufficiency} in the main text.
\end{definition}

\begin{definition}[Merge-consistency condition (C2)]
For every pair of strings $u,v\in\Strings$,
\[
\E\Big[\metric\Big(\fstar(u\concat v),\,\fstar\big(g(g(u)\concat g(v))\big)\Big)\Big]=0.
\]
The triangle inequality together with Condition~(C1) immediately implies $\E[\metric(\fstar(g(u)\concat g(v)),\,\fstar(g(g(u)\concat g(v))))]=0$, which is the specific check enforced at realized internal nodes in the hierarchy (Eq.~\eqref{eq:leaf_compatibility}).
\end{definition}

\begin{definition}[Idempotence condition (C3)]
For any random variable $Z$ supported on the range of $g$,
\[
\E\big[\metric\big(\fstar(g(Z)),\fstar(Z)\big)\,\big|\,Z\big]=0,
\]
ensuring that re-summarizing a summary leaves the oracle unchanged (Eq.~\eqref{eq:leaf_idempotence}).
\end{definition}

The edge-wise two-route comparison
\[
\metric\Big(\fstar(v\concat w),\fstar\big(g(v\concat w)\big)\Big)=0
\]
for realized child summaries $(v,w)$ is an optional audit probe (Remark~\ref{rem:edgewise_compatibility}).

\begin{remark}[Task-level congruence]
Whenever the distortion is zero, we may declare $x\sim x'$ if $\metric(\fstar(x),\fstar(x'))=0$. This equivalence relation is a congruence on the free monoid (Appendix~\ref{sec:global}) and motivates the view that $g$ should output canonical representatives of task-equivalent strings.
\end{remark}

\paragraph{Notation snapshot.}
\subsection{Nodewise preservation and the first theorem}

\begin{lemma}[Nodewise preservation under the consistency conditions]\label{lem:nodewise}
If Conditions~\eqref{eq:leaf_sufficiency} and \eqref{eq:leaf_compatibility} hold on all realized leaves and internal nodes of $T$, then for every node $u$,
\[
\E\Big[\metric\big(\fstar(g(u)),\fstar(S(u))\big)\Big]=0.
\]
\end{lemma}

\begin{proof}
If $u$ is a leaf, $g(u)=g(S(u))$ and Condition~(C1) yields the claim. If $u$ is internal with children $u_L,u_R$, the induction hypothesis gives $g(u_L)\sim S(u_L)$ and $g(u_R)\sim S(u_R)$. Condition~(C2) applied to the raw inputs $S(u_L),S(u_R)$ implies
\[
S(u_L)\concat S(u_R)\ \sim\ g\big(g(S(u_L))\concat g(S(u_R))\big).
\]
Because $g(S(u_L))=g(u_L)$ and $g(S(u_R))=g(u_R)$ by definition, the right-hand side equals $g(g(u_L)\concat g(u_R))=g(u)$. Thus $g(u)\sim S(u_L)\concat S(u_R)=S(u)$, completing the induction.
\end{proof}

\begin{proof}[Proof of Theorem~\ref{thm:one-pass}]
\phantomsection\label{proof:thm-one-pass}
Apply Lemma~\ref{lem:nodewise} to the root $r$ of $T$. The conditional expectation $\E_r$ coincides with $\E$ because every coin flip used anywhere in the reduction of $T$ sits below $r$. Moreover $S(r)=x$, establishing the claim.
\end{proof}

This lemma–theorem pair captures everything needed before invoking more elaborate schedules (Corollary~\ref{cor:schedule}) or additional rounds. From this point forward we simply write $g(u)$ for the summary stored at node $u$.

\subsection{Formal Development Cross-Reference}

The main results in this paper have been mechanically verified in Lean~4. Table~\ref{tab:lean-xref} provides a mapping between paper statements and their formal counterparts.

\begin{table}[h]
\centering
\small
\begin{tabular}{lll}
\toprule
Paper & Lean Theorem/Definition & File \\
\midrule
\multicolumn{3}{l}{\emph{Conditions and Definitions}} \\
Condition (C1) & \texttt{L1} (alias \texttt{C1}) & \texttt{LocalLaws.lean} \\
Condition (C2) & \texttt{L2} (alias \texttt{C2}) & \texttt{LocalLaws.lean} \\
Condition (C3) & \texttt{L3} (alias \texttt{C3}) & \texttt{LocalLaws.lean} \\
Global (A1)--(A3) & \texttt{A1\_global}, \texttt{A2\_global}, \texttt{A3\_global} & \texttt{GlobalAssumptions.lean} \\
\midrule
\multicolumn{3}{l}{\emph{Main Theorems}} \\
Theorem~\ref{thm:one-pass} & \texttt{one\_pass} & \texttt{PreservationTheorems.lean} \\
Theorem~\ref{thm:multi-round} & \texttt{multi\_round} & \texttt{ExpectationTheory.lean} \\
Corollary~\ref{cor:schedule} & \texttt{schedule\_invariance} & \texttt{PreservationTheorems.lean} \\
Corollary~\ref{cor:folds} & \texttt{fold\_of\_folds} & \texttt{PreservationTheorems.lean} \\
\midrule
\multicolumn{3}{l}{\emph{DPO Results (Appendix~\ref{sec:dpo})}} \\
Theorem~\ref{thm:dpo-exact} & \texttt{dpo\_exact\_metric}, \texttt{dpo\_equivalence} & \texttt{DPO.lean} \\
Theorem~\ref{thm:dpo-gap} & \texttt{dpo\_gap}, \texttt{dpo\_gap\_reward} & \texttt{DPO.lean} \\
\midrule
\multicolumn{3}{l}{\emph{Score Transport and Audit (Appendix~\ref{app:blackwell}--\ref{app:practical_audit})}} \\
Proposition~4 & \texttt{prop4\_cf\_implies\_oracle\_measurable} & \texttt{ScoreTransport.lean} \\
Proposition~5 & \texttt{prop5\_score\_transport} & \texttt{ScoreTransport.lean} \\
Proposition~7 & \texttt{prop7\_audit\_bound} & \texttt{AuditBounds.lean} \\
\midrule
\multicolumn{3}{l}{\emph{Existence and Counterexamples (Appendix~\ref{app:existence}, \ref{sec:examples})}} \\
Proposition~9.1 & \texttt{prop9\_1\_canonical\_rep\_exists} & \texttt{CounterexampleExistence.lean} \\
Proposition~9.2 & \texttt{prop9\_2\_compact\_encoding\_exists} & \texttt{CounterexampleExistence.lean} \\
Theorem~10.1 & \texttt{thm10\_1\_L3\_not\_derivable} & \texttt{CounterexampleExistence.lean} \\
\bottomrule
\end{tabular}
\caption{Cross-reference between paper statements and formal Lean development in \texttt{FormalProofs/}.}
\label{tab:lean-xref}
\end{table}

\noindent\textbf{Note on L2 vs.\ C2:} The formal L2 directly checks the composed condition $u\concat v \sim g(g(u)\concat g(v))$, which is equivalent to the conjunction of both parts of Condition~(C2) by transitivity. The two-part decomposition in the paper remains useful for auditing.

\subsection{Paper--Lean--Code Terminology Mapping}

Table~\ref{tab:terminology} provides a complete mapping between notation used in the paper, the Lean formalization, and the Python implementation.

\begin{table}[h]
\centering
\small
\begin{tabular}{lll}
\toprule
Paper & Lean & Python Code \\
\midrule
\multicolumn{3}{l}{\emph{Local Conditions}} \\
C1 (Sufficiency) & L1, C1 & \texttt{check\_type="sufficiency"} \\
C2 (Idempotence) & L3, C3 & \texttt{check\_type="idempotence"} \\
C3 (Merge consistency) & L2, C2 & \texttt{check\_type="merge\_consistency"} \\
--- & --- & \texttt{check\_type="substitution"} (C3 Case A) \\
\midrule
\multicolumn{3}{l}{\emph{Global Axioms}} \\
Assumption~1 & A1\_global & (not implemented) \\
Assumption~2 & A2\_global & (not implemented) \\
Assumption~3 & A3\_global & (not implemented) \\
\midrule
\multicolumn{3}{l}{\emph{Core Functions}} \\
$g : \Strings \rightsquigarrow \Strings$ & \texttt{g : α → PMF α} & \texttt{SummarizationStrategy.summarize()} \\
$\fstar : \Strings \to \Y$ & \texttt{fstar : α → Y} & \texttt{ScoringOracle.score()} \\
$\metric : \Y \times \Y \to [0,\infty)$ & \texttt{D : Y → Y → ℝ} & \texttt{1.0 - OracleScore.score} \\
$u \concat v$ & \texttt{u * v} & \texttt{format\_merge\_input(u, v)} \\
\midrule
\multicolumn{3}{l}{\emph{Rates and Bounds}} \\
$p_{\mathrm{suff}}$ & \texttt{pLeafAvg} & \texttt{AuditReport.sufficiency\_rate} \\
$p_{\mathrm{assoc}}$ & \texttt{pMergeAvg} & \texttt{AuditReport.assoc\_rate} \\
$p_{\mathrm{idem}}$ & \texttt{pIdemp} & \texttt{AuditReport.idempotence\_rate} \\
$\Delta_R$ & \texttt{Exp[D(ZR, x)]} & \texttt{get\_probabilistic\_bound()} \\
\midrule
\multicolumn{3}{l}{\emph{Guarantee Levels}} \\
Exact (gap = 0) & \texttt{EXACT} & \texttt{GuaranteeLevel.EXACT} \\
Union bound & \texttt{UNION\_BOUND} & \texttt{GuaranteeLevel.UNION\_BOUND} \\
Empirical & \texttt{EMPIRICAL} & \texttt{GuaranteeLevel.EMPIRICAL} \\
\midrule
\multicolumn{3}{l}{\emph{Sample Complexity}} \\
$\sqrt{\frac{\ln(2/\delta)}{2n}}$ & \texttt{confidence\_margin} & \texttt{confidence\_margin(delta, n)} \\
$\frac{\ln(2/\delta)}{2\epsilon^2}$ & \texttt{sample\_complexity} & \texttt{sample\_complexity(epsilon, delta)} \\
\bottomrule
\end{tabular}
\caption{Terminology mapping between paper notation, Lean formalization, and Python implementation. Note that the paper's C2/C3 numbering differs from Lean's L2/L3 ordering.}
\label{tab:terminology}
\end{table}

\noindent\textbf{Important notation difference:} The paper uses C2 for idempotence and C3 for merge consistency, while the Lean formalization uses L2 for merge consistency and L3 for idempotence. The Lean code provides \texttt{C1}, \texttt{C2}, \texttt{C3} as aliases that match the paper notation.

\section{Proofs of Propositions and Theorems in the Main Text}
\label{app:proofs-core}

This section collects the proofs of all corollaries and theorems stated in Section~\ref{sec:setup}. These results establish three key invariance properties that make hierarchical processing practical: (i)~schedule invariance—the expected oracle is the same whether we use balanced trees or daisy chains; (ii)~fold-of-folds invariance—we can nest hierarchical reductions arbitrarily; and (iii)~multi-round preservation—repeatedly ``normalizing'' summaries does not drift the oracle when on-range stability holds. All three follow from the consistency conditions (C1)--(C3) via structural induction.

Throughout this section we fix a document $x=b_1\concat\cdots\concat b_k\in\Strings$ together with a realized full binary reduction tree $T$ on the leaves $(b_1,\dots,b_k)$. For a node $u\in T$ we write $S(u)$ for the realized string at $u$ and $g(u)$ for the corresponding hierarchical evaluation defined in Appendix~\ref{app:notation}. Expectations with $\E$ average \emph{only} over the internal randomness of $g$ used in the computation of $g(u_L)$, $g(u_R)$, and the parent call at $u$, conditional on the realized subtree below~$u$.

\medskip\noindent
Every result below works solely with the task pseudometric $\metric$; equalities are always stated after applying $\fstar$. Lemma~\ref{lem:nodewise} and Theorem~\ref{thm:one-pass} already establish nodewise and root preservation under the three consistency conditions. Here we spell out three consequences: schedule invariance (any binary parenthesization on a fixed partition yields the same expected oracle), invariance under ``fold-of-folds'' execution plans, and multi-round preservation provided $g$ is stable on its own range.

\begin{proof}[Proof of Corollary~\ref{cor:schedule} (Schedule invariance)]
\phantomsection\label{proof:cor-schedule}
Lemma~\ref{lem:nodewise} applies to \emph{any} realized node set. Changing the parenthesization merely changes which internal nodes appear, so the lemma re-runs verbatim and Theorem~\ref{thm:one-pass} yields the same root equality.
\end{proof}


\begin{proof}[Proof of Corollary~\ref{cor:folds} (Fold-of-folds invariance)]
\phantomsection\label{proof:cor-folds}
Let $(F_1,\ldots,F_J)$ be a contiguous coarsening of the leaves. For each $j$, fix an arbitrary full
binary within‑fold tree $T_j$ whose leaves, read left to right, concatenate to $F_j$. Let $\widehat T$ be any full
binary over‑fold tree whose leaves, read left to right, concatenate to $F_1+\cdots+F_J$. Form the
\emph{composite tree} $T_{\mathrm{comp}}$ with root $r$ by grafting $T_1,\ldots,T_J$ into $\widehat T$ at its $J$ leaves.
If \eqref{eq:leaf_sufficiency} holds on the realized leaves of $T_{\mathrm{comp}}$ and \eqref{eq:leaf_compatibility} holds at every realized internal node of $T_{\mathrm{comp}}$, then
\[
\E\!\Big[\,\metric\!\big(\fstar(g(r)),\ \fstar(x)\big)\,\Big]\;=\;0.
\]

In particular, the expected oracle at the root is invariant to the choice of the within-fold parenthesizations and to the over-fold parenthesization, provided the realized leaves and internal nodes of $T_{\mathrm{comp}}$ satisfy \eqref{eq:leaf_sufficiency}–\eqref{eq:leaf_compatibility}. Moreover, if \eqref{eq:leaf_idempotence} (on-range stability) also holds, then for the usual two-level “fold-of-folds” computation that first reduces each $F_j$ to $S_j:=g(T_j)$ and then reduces $(S_1,\ldots,S_J)$ along $\widehat T$, we have
\[
\E\!\Big[\,\metric\!\big(\fstar(g(\widehat T)),\ \fstar(x)\big)\,\Big]\;=\;0,
\]
because all additional $g$ calls at the leaves of $\widehat T$ are re‑applications of $g$ to on‑range strings and are inert at the oracle by \eqref{eq:leaf_idempotence}.
Build the composite tree $T_{\mathrm{comp}}$ that first runs $T_1,\ldots,T_J$ and then grafts their outputs into $\widehat T$. Every realized leaf or internal node across both levels appears exactly once inside $T_{\mathrm{comp}}$, so Lemma~\ref{lem:nodewise} and Theorem~\ref{thm:one-pass} applied to $T_{\mathrm{comp}}$ give the displayed equality. For the two-stage implementation, note that the leaves of $\widehat T$ are already on the range of $g$, so on-range stability \eqref{eq:leaf_idempotence} makes the extra applications of $g$ inert in expectation.
\end{proof}


% \begin{corollary}[Fold‑of‑folds invariance; Corollary~\ref{cor:folds}]
% \label{cor:folds-appendix}
% Let $(F_1,\ldots,F_J)$ be a contiguous coarsening of the leaves. Reduce each fold $F_j$ by an arbitrary within‑fold tree $T_j$ to obtain $S_j:=g(T_j)$, then reduce $(S_1,\ldots,S_J)$ by any over‑fold tree $\widehat T$. If \eqref{eq:leaf_sufficiency} holds on all realized leaves across both levels and \eqref{eq:leaf_compatibility} holds at every realized internal node, then
% \[
% \E\!\big[\ \metric\big(\fstar\big(g(\widehat T)\big),\ \fstar(x)\big)\big]\;=\;0.
% \]
% \end{corollary}

% \begin{proof}
% Fix $j$ and condition on the $\sigma$‑field below $\mathrm{root}(T_j)$. Because \eqref{eq:leaf_sufficiency} holds on the realized leaves of $T_j$ and \eqref{eq:leaf_compatibility} holds at every realized internal node of $T_j$, Lemma~\ref{lem:nodewise} yields $\E_{\mathrm{root}(T_j)}\!\big[\metric\big(\fstar(S_j),\fstar(F_j)\big)\big]=0$. Now view $\widehat T$ as a tree with leaves $(S_1,\ldots,S_J)$. Since \eqref{eq:leaf_sufficiency} also holds on these realized leaves and \eqref{eq:leaf_compatibility} holds at every realized internal node of $\widehat T$, a second application of Lemma~\ref{lem:nodewise} gives
% \[
% \E_{\mathrm{root}(\widehat T)}\!\Big[\ \metric\!\big(\fstar\big(g(\widehat T)\big),\ \fstar(S_1\concat\cdots\concat S_J)\big)\ \Big]\;=\;0.
% \]
% By \eqref{eq:leaf_sufficiency}, replacing each $F_j$ by $S_j$ leaves the oracle target unchanged, so $\fstar(S_1\concat\cdots\concat S_J)=\fstar(F_1\concat\cdots\concat F_J)=\fstar(x)$. Taking expectations and using the tower property, the displayed conditional equality lifts to $\E\!\big[\metric\big(\fstar(g(\widehat T)),\fstar(x)\big)\big]=0$, as claimed.
% \end{proof}


\begin{proof}[Proof of Theorem~\ref{thm:multi-round} (Multi-round preservation)]
\phantomsection\label{proof:thm-multi-round}
Set $\Delta_R(x):=\E[\metric(\fstar(Z^{(R)}(x)),\fstar(x))]$. The base case $\Delta_1(x)=0$ is Theorem~\ref{thm:one-pass}. For the inductive step, use stability on the range of $g$ to note that $\E[\metric(\fstar(g(Z^{(R)}(x))),\fstar(Z^{(R)}(x)))] = 0$. The triangle inequality thus gives $\Delta_{R+1}(x)\le \Delta_R(x)$, and the sequence remains zero.
\end{proof}

\input{appendix/03_blackwell_score_transport.tex}
\section{Practical audit guidance}\label{app:practical_audit}

What matters at deployment is a single number,
\[
\Delta_R\ :=\ \E\Big[\,\metric\big(\fstar(Z^{(R)}(X)),\ \fstar(X)\big)\,\Big],
\]
the expected \emph{task–space} distortion after one hierarchical pass and $R-1$ optional
re‑summarization rounds. All statements below assume $\metric\in[0,1]$ (rescale if needed by
$\mathrm{diam}(Y)$). The audit is local, deployment‑specific, and piggybacks exactly on the two‑route
merge diagram: a parent can be reached either by “concatenate–summarize’’ or by “summarize
children–concatenate–summarize again’’ (Figure~\ref{fig:local_compatibility}).%

\paragraph{Exact pass short‑circuit.}
If the realized leaves satisfy \textnormal{(C1)} exactly and the realized internal nodes satisfy
\textnormal{(C2)} exactly, Theorem~\ref{thm:one-pass} gives $\Delta_1=0$. If, in addition, \textnormal{(C3)}
holds, then Theorem~\ref{thm:multi-round} yields $\Delta_R=0$ for all $R\ge2$. In this regime there is
nothing to bound; the audit reduces to checking that the empirical versions of the three quantities
below are (within tolerance) zero.

\paragraph{What to estimate (local violation rates).}
Let a realized reduction tree have $N$ leaves and $M=N-1$ internal nodes. Denote realized
leaves by $b$ (each $b$ is a contiguous substring of the base document $x_0$ and is never itself a summary), and realized internal nodes by $u$ with children $u_L,u_R$ whose stored summaries are $g(u_L),g(u_R)$. Define the \emph{indicator} violations
\[
p_{\text{suff}}\ :=\ \E\,\mathbf{1}\!\left\{\metric\!\big(\fstar(g(B)),\,\fstar(B)\big)>0\right\},
\]
\[
p_{\text{bound}}\ :=\ \E\,\mathbf{1}\!\left\{\metric\!\big(\fstar(g(B_i\concat B_{i+1})),\,\fstar(g(g(B_i)\concat g(B_{i+1})))\big)>0\right\},
\]
\[
p_{\text{merge}}\ :=\ \E\,\mathbf{1}\!\left\{\metric\!\big(\fstar(S_L\concat S_R),\,\fstar(g(S_L\concat S_R))\big)>0\right\},
\]
\[
p_{\text{idem}}\ :=\ \E\,\mathbf{1}\!\left\{\metric\!\big(\fstar(g(Z^{(1)}(x))),\,\fstar(Z^{(1)}(x))\big)>0\right\}\qquad\text{(only if $R\ge2$).}
\]
Let $\lambda\in[0,1]$ denote the probability that a merge-consistency sample targets a leaf boundary (so $1-\lambda$ is the probability it targets an internal merge). The aggregate violation rate used by the bounds below is
\[
p_{\text{assoc}}\ :=\ \lambda\,p_{\text{bound}} + (1-\lambda)\,p_{\text{merge}}.
\]
The two components correspond exactly to the two links in Condition~\eqref{eq:leaf_compatibility}: $p_{\text{bound}}$ measures the joint-vs-disjoint comparison when raw concatenations fit in context, while $p_{\text{merge}}$ measures the raw-vs-summary comparison on stored child strings.
In code: log each raw leaf $b$ (the substring of $x_0$ that entered the hierarchy), its summary $g(b)$, and, for every internal node, the ordered pair of stored child strings $\big(g(u_L),g(u_R)\big)$ together with the resulting parent summary $g(u)$. Also record the seeds/call counts used by $g$ so runs can be replayed. Then \\emph{sample} leaves and internal nodes, \\emph{replay} the same randomness policy when needed, evaluate the displayed $\metric$-terms via $\fstar$, threshold at $>0$, and average. Conditioning on the realized subtree in the merge term matches \eqref{eq:leaf_compatibility} exactly.

\paragraph{Crude but rigorous aggregation by union bound.}
Define the event that the root deviates after $R$ rounds by
\[
\mathcal{E}_R\ :=\ \left\{\ \metric\!\big(\fstar(Z^{(R)}(X)),\ \fstar(X)\big)>0\ \right\}.
\]
If no leaf fails Condition~\eqref{eq:leaf_sufficiency} and no realized internal node fails Condition~\eqref{eq:leaf_compatibility} (and, for $R\ge2$, no on‑range
re‑summarization fails Condition~\eqref{eq:leaf_idempotence} at the root between rounds), then Theorems~\ref{thm:one-pass}
and~\ref{thm:multi-round} force $\mathcal{E}_R$ not to occur. Therefore

% \[
% \mathcal{E}_R\ \subseteq\ \bigcup_{\text{leaves }b}\left\{\metric(\fstar(g(b)),\fstar(b))>0\right\}
% \ \cup\ \bigcup_{\text{internal }u}\left\{\metric\!\Big(\fstar\!\big(g(g(u_L)\!+\!g(u_R)\big),\ \fstar\!\big(g(u_L)\!+\!g(u_R)\big)\Big)>0\right\}
% \ \cup\ \bigcup_{t=1}^{R-1}\left\{\metric\!\big(\fstar(g(Z^{(t)}(X))),\fstar(Z^{(t)}(X))\big)>0\right\}.
% \]

\begin{align*}
\mathcal{E}_R \subseteq{}& \bigcup_{\text{leaves }b}\left\{\metric(\fstar(g(b)),\fstar(b))>0\right\} \\
&\cup \bigcup_{\text{internal }u}\left\{\metric\!\Big(\fstar\!\big(g(g(u_L)\!+\!g(u_R)\big),\ \fstar\!\big(g(u_L)\!+\!g(u_R)\big)\Big)>0\right\} \\
&\cup \bigcup_{t=1}^{R-1}\left\{\metric\!\big(\fstar(g(Z^{(t)}(X))),\fstar(Z^{(t)}(X))\big)>0\right\}.
\end{align*}


Taking probabilities and applying the union bound gives the \emph{crude envelope}
% \[
% \Pr(\mathcal{E}_1)\ \le\ N\,p_{\text{suff}}\ +\ M\,p_{\text{assoc}},
% \qquad
% \Pr(\mathcal{E}_R)\ \le\ N\,p_{\text{suff}}\ +\ M\,p_{\text{assoc}}\ +\ (R-1)\,p_{\text{idem}}\quad(R\ge2).
% \]

\begin{align*}
\Pr(\mathcal{E}_1) &\le N\,p_{\text{suff}} + M\,p_{\text{assoc}}, \\
\Pr(\mathcal{E}_R) &\le N\,p_{\text{suff}} + M\,p_{\text{assoc}} + (R-1)\,p_{\text{idem}}\quad(R\ge2).
\end{align*}

Since $\metric\in[0,1]$, we have $\Delta_R=\E[\metric(\cdot,\cdot)]\le \Pr(\mathcal{E}_R)$, hence
\begin{equation}\label{eq:unionbound}
\Delta_1\ \le\ N\,p_{\text{suff}}\ +\ M\,p_{\text{assoc}},\qquad
\Delta_R\ \le\ N\,p_{\text{suff}}\ +\ M\,p_{\text{assoc}}\ +\ (R-1)\,p_{\text{idem}}\quad(R\ge2).
\end{equation}
These bounds hold with no assumption beyond the consistency conditions used by Theorems~\ref{thm:one-pass}
and~\ref{thm:multi-round}. They are intentionally coarse: any single local failure is allowed to
account for a root deviation. When \eqref{eq:leaf_idempotence} holds, $p_{\text{idem}}=0$ and the
multi‑round term vanishes.

\paragraph{Plug‑in estimates.}
In deployment we report the sample plug‑ins

\begin{align*}
\widehat p_{\text{suff}} &:= \frac{1}{n_\ell}\sum_{i=1}^{n_\ell}\mathbf{1}\!\left\{\metric\!\big(\fstar(g(b_i)),\fstar(b_i)\big)>0\right\}, \\
\widehat p_{\text{bound}} &:= \frac{1}{n_{\text{bound}}}\sum_{j=1}^{n_{\text{bound}}}\mathbf{1}\!\left\{\metric\!\Big(\fstar\!\big(g(b_{j}\concat b_{j+1})\big),\ \fstar\!\big(g(g(b_{j})\concat g(b_{j+1}))\big)\Big)>0\right\}, \\
\widehat p_{\text{merge}} &:= \frac{1}{n_{\text{merge}}}\sum_{j=1}^{n_{\text{merge}}}\mathbf{1}\!\left\{\metric\!\Big(\fstar\!\big(s_{L,j}\concat s_{R,j}\big),\ \fstar\!\big(g(s_{L,j}\concat s_{R,j})\big)\Big)>0\right\},
\end{align*}
and, if $R\ge2$,
\begin{align*}
\widehat p_{\text{idem}} &:= \frac{1}{n_r}\sum_{k=1}^{n_r}\mathbf{1}\!\left\{\metric\!\big(\fstar(g(Z^{(1)}(x_k))),\,\fstar(Z^{(1)}(x_k))\big)>0\right\}.
\end{align*}
Setting $\widehat\lambda:=\frac{n_{\text{bound}}}{n_{\text{bound}}+n_{\text{merge}}}$ yields
\[
\widehat p_{\text{assoc}}:=\widehat\lambda\,\widehat p_{\text{bound}} + (1-\widehat\lambda)\,\widehat p_{\text{merge}}.
\]
The \emph{crude deployment-level bound} is then
\begin{align*}
\widehat\Delta_1^{\,\mathrm{ub}} &:= N\,\widehat p_{\text{suff}} + M\,\widehat p_{\text{assoc}}, \\
\widehat\Delta_R^{\,\mathrm{ub}} &:= N\,\widehat p_{\text{suff}} + M\,\widehat p_{\text{assoc}} + (R-1)\,\widehat p_{\text{idem}}\quad(R\ge2).
\end{align*}
% \[
% \widehat p_{\text{suff}}\ :=\ \frac1{n_\ell}\sum_{i=1}^{n_\ell}\mathbf{1}\!\left\{\metric\!\big(\fstar(g(b_i)),\fstar(b_i)\big)>0\right\},
% \quad
% \widehat p_{\text{assoc}}\ :=\ \frac1{n_m}\sum_{j=1}^{n_m}\mathbf{1}\!\left\{\metric\!\Big(\fstar\!\big(g(g(u_{L,j})\!+\!g(u_{R,j})\big),\ \fstar\!\big(g(u_{L,j})\!+\!g(u_{R,j})\big)\Big)>0\right\},
% \]
% and, if $R\ge2$,
% \[
% \widehat p_{\text{idem}}\ :=\ \frac1{n_r}\sum_{k=1}^{n_r}\mathbf{1}\!\left\{\metric\!\big(\fstar(g(Z^{(1)}(x_k))),\,\fstar(Z^{(1)}(x_k))\big)>0\right\}.
% \]
% The \emph{crude deployment‑level bound} is then
% \[
% \widehat\Delta_1^{\,\mathrm{ub}}\ :=\ N\,\widehat p_{\text{suff}}\ +\ M\,\widehat p_{\text{assoc}},
% \qquad
% \widehat\Delta_R^{\,\mathrm{ub}}\ :=\ N\,\widehat p_{\text{suff}}\ +\ M\,\widehat p_{\text{assoc}}\ +\ (R-1)\,\widehat p_{\text{idem}}\quad(R\ge2),
% \]
truncated at $1$ if desired. Sampling should be random over the realized leaves and realized
internal nodes of the run; if worried about depth or heterogeneity, stratify by depth or by child
summary length and then aggregate. Labeling is whatever $\fstar$ requires (scripted or human).

\paragraph{Direct root estimation (often simplest).}
When you do not need a decomposition, estimate $\Delta_R$ directly by sampling documents
$x_i$ and computing $\widehat\Delta_R=\frac1n\sum_i \metric(\fstar(Z^{(R)}(x_i)),\fstar(x_i))$.
If $R\ge2$ and \eqref{eq:leaf_idempotence} holds, $\Delta_R=\Delta_1$, so it suffices to measure $R=1$.

% \section{Practical audit guidance}\label{app:practical_audit}

% What matters at deployment is a single number,
% \[
% \Delta_R\ :=\ \E\Big[\,\metric\big(\fstar(Z^{(R)}(X)),\ \fstar(X)\big)\,\Big],
% \]
% the expected \emph{task–space} distortion after one hierarchical pass and $R-1$ optional re‑summarization rounds. We estimate $\Delta_R$ from the very operations the hierarchy executed: realized leaves, realized internal merges, and (only if we re‑summarize) an on‑range probe at the realized root. The audit is therefore local, deployment‑specific, and reproducible. It piggybacks exactly on the two‑route merge picture: a parent can be reached either by “concatenate–summarize’’ or by “summarize children–concatenate–summarize again’’ (Figure~\ref{fig:local_compatibility}).%

% \paragraph{Exact pass short‑circuit.}
% If the realized leaves satisfy \eqref{eq:leaf_sufficiency}exactly and the realized internal nodes satisfy \eqref{eq:leaf_compatibility} exactly, then Theorem~\ref{thm:one-pass} gives
% \[
% \E\,\metric\!\Big(\fstar\big(g(\mathrm{root}(T))\big),\ \fstar(x)\Big)\;=\;0,
% \]
% so $\Delta_1=0$. If, in addition, \eqref{eq:leaf_idempotence} holds, then Theorem~\ref{thm:multi-round} yields $\Delta_R=0$ for every $R\ge2$. In this exact regime there is nothing to bound; the audit reduces to checking that the empirical versions of the three local quantities below are (within tolerance) zero on realized leaves, realized merges, and on‑range re‑summaries.%

% \paragraph{What to estimate (operation‑level probes).}
% Let a realized reduction tree have $N$ leaves and $M=N-1$ internal nodes. Denote realized leaves by $b$, and realized internal nodes by $u$ with children $u_L,u_R$ and stored summaries $g(u_L),g(u_R)$. Using $\E$ for the internal randomness of $g$ and $\E$ for the conditional expectation over all coin flips used below $u$ and at its parent call, we define three observable quantities and estimate each by simple random sampling over the realized nodes:
% \[
% \delta_{\mathrm{leaf}}\ :=\ \E\,\metric\!\big(\fstar(g(B)),\,\fstar(B)\big),
% \]
% \[
% \epsilon_{\mathrm{merge}}\ :=\ \E\ \metric\!\Big(\fstar\!\big(g(g(u_L)\concat g(u_R))\big),\ \fstar\!\big(g(u_L)\concat g(u_R)\big)\Big),
% \]
% \[
% \iota_{\mathrm{root}}\ :=\ \E\,\metric\!\big(\fstar(g(Z^{(1)}(x))),\,\fstar(Z^{(1)}(x))\big)\qquad\text{(only if $R\ge2$).}
% \]
% These target, respectively, leaf sufficiency \eqref{eq:leaf_sufficiency}, parentwise compatibility at realized merges \eqref{eq:leaf_compatibility}, and on‑range stability for extra rounds \eqref{eq:leaf_idempotence}. In code: log $S(\cdot)$ and $g(\cdot)$ at every realized node together with the seeds/call‑counts used by $g$; then \emph{sample} leaves and internal nodes, \emph{replay} the same randomness policy when needed, evaluate the displayed $\metric$‑terms via $\fstar$, and average. Conditioning on the realized subtree in the merge term keeps the estimate aligned with the nodewise statement of \eqref{eq:leaf_compatibility}.%


% \paragraph{Aggregation.}
% There are two clean ways to turn the node‑level probes into a deployment‑level number; both are faithful to the realized schedule.

% \noindent\emph{Option A: direct root estimation (no further assumptions).}
% The most straightforward approach is directly aggregating at the root level. For any fixed $R\ge1$, draw a random sample of documents $x_i$ from the deployment and compute
% \[
% \widehat\Delta_R\ :=\ \frac{1}{n}\sum_{i=1}^n \metric\!\big(\fstar(Z^{(R)}(x_i)),\,\fstar(x_i)\big).
% \]
% This is an unbiased estimator of $\Delta_R$. When $R\ge2$, \eqref{eq:leaf_idempotence} implies $Z^{(R)}$ is on‑range and additional rounds are inert in expectation (Theorem~\ref{thm:multi-round}), so the direct estimate for $R=1$ already controls all further rounds.

% \noindent\emph{Option B: additive envelope under a testable parent‑stability law.}
% If a decomposition over realized leaves and internal merges is desired, introduce the following local, on‑range hypothesis at the level of parent (summary) nodes:
% \begin{assumption}[On‑range determinacy and non‑expansiveness at parents]\label{ass:det-lip}
% There exists a measurable $M:\Y\times\Y\to\Y$ such that for all on‑range $z_L,z_R$,
% \[
% \fstar\!\big(g(z_L\concat z_R)\big)\ =\ M\big(\fstar(z_L),\,\fstar(z_R)\big)\quad\text{almost surely},
% \]
% and $M$ is $1$‑Lipschitz in each coordinate:
% \[
% \metric\!\big(M(y_1,y_2),\,M(y_1',y_2)\big)\ \le\ \metric(y_1,y_1'),\qquad
% \metric\!\big(M(y_1,y_2),\,M(y_1,y_2')\big)\ \le\ \metric(y_2,y_2').
% \]
% \end{assumption}
% Assumption~\ref{ass:det-lip} is strictly local to realized, on‑range inputs and is directly testable from deployment logs by holding child oracles fixed while varying their concrete realizations, and by perturbing one child-node's oracle output at a time. Under \eqref{ass:det-lip} one obtains the nodewise recursion
% \[
% D(u)\ :=\ \E\,\metric\!\big(\fstar(g(u)),\,\fstar(S(u))\big)\ \le\ D(u_L)\ +\ D(u_R)\ +\ \epsilon(u),
% \]
% with $D(b)=\delta_{\mathrm{leaf}}(b)$ at leaves and $\epsilon(u)$ as above at internal nodes. Summing up the realized tree yields the additive envelope for one pass,
% \[
% \Delta_1\ \le\ \sum_{\text{leaves }b}\delta_{\mathrm{leaf}}(b)\ +\ \sum_{\text{internal }u}\epsilon(u),
% \]
% and, for $R\ge2$,
% \[
% \Delta_R\ \le\ \Delta_1\ +\ (R-1)\,\iota_{\mathrm{root}}.
% \]
% When \eqref{eq:leaf_idempotence} holds, $\iota_{\mathrm{root}}=0$ and extra rounds are inert in expectation.%
 
% \paragraph{What to report.}
% Report $(N,M,R)$ for the realized run; report the sampled means $(\hat\delta_{\mathrm{leaf}},\,\hat\epsilon_{\mathrm{merge}},\,\hat\iota_{\mathrm{root}})$ together with either the direct root estimate $\widehat\Delta_R$.
% (Option~A) or the additive plug‑in summary $\,\widehat\Delta_1:=\sum_b\hat\delta_{\mathrm{leaf}}(b)+\sum_u\hat\epsilon_{\mathrm{merge}}(u)$ and $\widehat\Delta_R:=\widehat\Delta_1+(R-1)\hat\iota_{\mathrm{root}}$ under Assumption~\ref{ass:det-lip} (Option~B). 
% For uncertainty, add nonparametric bootstrap intervals over the sampled units. 
% % If $\Y=\{0,1\}$ and $\metric$ is Hamming, the same procedure returns flip shares; no new notation is needed.

% \paragraph{Why $g\in\Gstar$ does not subsume the parent test.}
% Exact $\fstar$‑sufficiency ($g\in\Gstar$) guarantees $\metric(\fstar(g(z)),\fstar(z))=0$ for every \emph{single} input $z$ and hence makes \eqref{eq:leaf_compatibility} vacuous and \eqref{eq:leaf_idempotence} hold whenever $g$ is idempotent on‑range. It does not constrain how $\fstar$ reacts to \emph{replacing both children by on‑range representatives and then concatenating}. That cross‑child behavior is precisely what \eqref{eq:leaf_compatibility} audits at realized merges. Thus, even with $g\in\Gstar$, you must still check \eqref{eq:leaf_compatibility} on the edges your scheduler actually realized. When \eqref{eq:leaf_sufficiency}–\eqref{eq:leaf_compatibility} (and \eqref{eq:leaf_idempotence} if $R\ge2$) hold exactly, the section’s “Exact pass” short‑circuit applies and yields $\Delta_R=0$; otherwise, use Option~A or, if you are willing to assume and test Assumption~\ref{ass:det-lip}, the additive Option~B.

\section{Local Consistency Bootstrap}\label{app:bootstrap}

This appendix instantiates the ``consistency bootstrap'' used throughout the main text. We view the hierarchical summarizer as a programmable policy $\mathcal{P}_\theta$ that exposes two callable modules:
\begin{itemize}
    \item \textbf{LeafSummarizer}$(b;\theta_\ell)$ maps a raw block $b$ to a summary $s_b$. Its loss is the indicator from Condition~\eqref{eq:leaf_sufficiency}.
    \item \textbf{MergeSummarizer}$((s_{u_L},s_{u_R});\theta_m)$ maps stored child summaries to a parent summary $s_u$. Its loss is the indicator from Condition~\eqref{eq:leaf_compatibility}.
\end{itemize}
DSPy \citep{khattab2024dspy} provides a natural implementation: each module is a DSPy ``program'' whose parameters are prompt slots, few-shot exemplars, and temperature controls. The bootstrap proceeds as follows.

\paragraph{1. Leaf pass.} Sample a batch of leaves $\{b_i\}$. For each leaf, run \textbf{LeafSummarizer}, evaluate $\metric(\fstar(g(b_i)),\fstar(b_i))$, and store the trace $(b_i, s_{b_i}, \text{oracle})$. When the oracle disagrees, regenerate $k$ candidates at elevated temperature or solicit a human rewrite. The first candidate that matches the oracle becomes a ``positive'' trace; if none exists, mark the edge as a hard negative for escalation. All traces are appended to DSPy's trace buffer and optionally distilled into few-shot exemplars.

\paragraph{2. Merge pass.} Using the best available leaf summaries, sample adjacent pairs $(s_{u_L},s_{u_R})$. Run \textbf{MergeSummarizer}, check $\metric(\fstar(g(s_{u_L}\concat s_{u_R})),\fstar(s_{u_L}\concat s_{u_R}))$, and record the trace. Successful merges are stored as exemplars. Failures trigger regeneration or escalation. Because the oracle only sees the concatenated child summaries, this step never requires reopening raw text.

\paragraph{3. Parameter update.} Run DSPy's optimizer (e.g., coordinate descent over prompt slots or gradient-free search over temperatures) on the accumulated traces. The objective is simply the empirical violation rate: minimize the fraction of traces where the oracle disagrees. Both modules can be optimized jointly or in alternating rounds. When operated online, we interleave sampling and updates so that fresh violations immediately produce corrective exemplars.

\paragraph{4. Stopping rule.} Maintain running estimates of $\hat{p}_{\text{suff}}$, $\hat{p}_{\text{assoc}}$, and (if applicable) $\hat{p}_{\text{idem}}$. Continue bootstrapping until these empirical rates fall below the deployment tolerance, as measured either on a dedicated validation set or via sequential confidence bounds. Hard negatives that never satisfy the oracle become the most valuable few-shot examples and should be preserved even after the rate target is met.

This loop requires no external supervision beyond the oracle assumed in the main body. Human edits or calls to a larger teacher model only appear when regeneration repeatedly fails, and those interventions are logged so that the improved spans become part of the next prompt. The resulting modules enforce the three consistency conditions by construction, aligning the entire hierarchy without ever inspecting the full document.

\section{Direct Preference Optimization (DPO): Equivalence and Gap Bounds}
\label{sec:dpo}

Direct preference optimization trains policies by comparing preferred (``win'') vs.\ dispreferred (``loss'') responses. This section shows when training on hierarchically summarized documents $Z^{(R)}(X)$ yields the same optimal policies as training on full documents $X$. Under exact oracle preservation—when Conditions~\eqref{eq:leaf_sufficiency}--\eqref{eq:leaf_idempotence} hold exactly along the realized tree—the DPO objective depends on the input only through $\fstar(\cdot)$, so the sets of population minimizers coincide. When preservation is approximate, we control the population gap by the oracle distortion $\Delta_R$ and provide explicit bounds under Lipschitz or bounded-reward assumptions. The key requirement is that annotator preferences factor through the oracle; if preferences depend on presentational features not encoded by $\fstar$, training on summaries can deviate even when oracle preservation holds.

Let preferences follow Bradley–Terry–Luce with per‑oracle rewards $\{R_y:\A\to\R\}_{y\in\Y}$ and scale $\beta>0$:
\[
\Pr\{a^w\succ a^\ell\mid X=x\}=\sigma\!\Big(\beta\,[R_{\fstar(x)}(a^w)-R_{\fstar(x)}(a^\ell)]\Big),
\qquad \sigma(t)=\frac{1}{1+e^{-t}}.
\]
DPO reparameterizes the reward through the policy $\pi$ relative to a reference policy $\pi_{\mathrm{ref}}$ and minimizes
\[
\mathcal L_{\mathrm{DPO}}(\pi;\pi_{\mathrm{ref}})=
-\,\E_{(X,a^w,a^\ell)}\!\left[\log \sigma\!\Big(\beta\Big\{\log\frac{\pi(a^w\mid X)}{\pi_{\mathrm{ref}}(a^w\mid X)}-\log\frac{\pi(a^\ell\mid X)}{\pi_{\mathrm{ref}}(a^\ell\mid X)}\Big\}\Big)\right].
\]
The log-ratio difference $\log\frac{\pi(a^w\mid X)}{\pi_{\mathrm{ref}}(a^w\mid X)}-\log\frac{\pi(a^\ell\mid X)}{\pi_{\mathrm{ref}}(a^\ell\mid X)}$ measures how much more the policy $\pi$ favors the winner over the loser, relative to the reference. DPO trains $\pi$ to maximize this gap for preferred pairs.

\begin{definition}[Oracle‑measurable policies]
A policy $\pi$ (and $\pi_{\mathrm{ref}}$) is oracle‑measurable if $\pi(\cdot\mid x)=\pi(\cdot\mid x')$ whenever $\metric\!\big(\fstar(x),\fstar(x')\big)=0$.
\end{definition}

\begin{theorem}[Exact DPO equivalence under oracle preservation and oracle–indexed pairs]\label{thm:dpo-exact}
% \emph{(Restates Theorem~\ref{thm:dpo-exact}.)}
Assume $\metric\!\big(\fstar(Z^{(R)}(X)),\fstar(X)\big)=0$ almost surely for some $R\ge1$. Assume also that the pair generator depends on $X$ only through $\fstar(X)$ and that $\pi_{\mathrm{ref}}$ is oracle‑measurable with full support. Then the sets of population minimizers of $\mathcal L_{\mathrm{DPO}}$ coincide when trained on $X$ and when trained on $Z^{(R)}(X)$. Moreover, there exists a population minimizer that is oracle‑measurable; equivalently, we may without loss of generality restrict attention to oracle‑measurable policies.
\end{theorem}

\begin{proof}
Write $Y=\fstar(X)$ and note that exact preservation implies $\fstar(Z^{(R)}(X))=Y$ almost surely. Under the oracle‑indexed pair generator and oracle‑measurable $\pi_{\mathrm{ref}}$, the joint law of $(Y,a^w,a^\ell)$ is the same whether the input to the learning objective is $X$ or $Z^{(R)}(X)$, and the two inner logits depend on the input only through $Y$. Therefore $\mathcal L_{\mathrm{DPO}}$ decomposes as an expectation over $Y$ of a functional that depends on $\pi(\cdot\mid\cdot)$ only through its $Y$–sections. For each $y$, any Bayes–optimal $\pi(\cdot\mid y)$ solves an identical one‑dimensional problem under $X$ and under $Z^{(R)}(X)$, so the sets of minimizers coincide and every minimizer is $Y$–measurable (oracle‑measurable).
\end{proof}

\begin{theorem}[Quantitative DPO gap under metric distortion (Sketch)]\label{thm:dpo-gap}
% \emph{(Restates Theorem~\ref{thm:dpo-gap}.)}
\textbf{Note:} This quantitative bound is a proof sketch. The formal Lean verification in \texttt{FormalProofs/DPO.lean} proves the \emph{exact} case (Theorem~\ref{thm:dpo-exact}); the Lipschitz bounds below are provided as intuition but are not machine-verified. See Remark~(iii) for what IS formally proven.

Let $\Delta_R=\E\big[\,\metric\!\big(\fstar(Z^{(R)}(X)),\fstar(X)\big)\,\big]$ be the expected oracle distortion.\footnote{This bound assumes the pair generator is either constant (independent of the document) or oracle-indexed. The formal Lean proof uses a constant generator for the tighter bound; see Remark~(iii) below for details.} To bound the population gap quantitatively, we impose regularity conditions ensuring that policies do not react radically to small changes in the oracle value. One option is a \emph{policy-Lipschitz log-ratio} condition: there exists $L_\pi<\infty$ such that
\[
\Bigg|\log\frac{\pi(a\mid y)}{\pi_{\mathrm{ref}}(a\mid y)}-\log\frac{\pi(a\mid y')}{\pi_{\mathrm{ref}}(a\mid y')}\Bigg|\,\le\,L_\pi\,\metric(y,y')\quad\forall\,a,y,y',
\]
so the policy's preference for action $a$ changes smoothly with the oracle value. Alternatively, we may assume a \emph{reward-Lipschitz exponential family}: $\pi(a\mid y)\propto \pi_{\mathrm{ref}}(a\mid y)\exp\{\beta^{-1}R_y(a)\}$ with $\sup_a|R_y(a)-R_{y'}(a)|\le L_R\,\metric(y,y')$ for all $y,y'$, which says the induced rewards vary smoothly in the oracle space.
Then, writing $\mathcal L^{X}$ and $\mathcal L^{Z}$ for the DPO loss with inputs $X$ and $Z^{(R)}(X)$, respectively,
\[
\big|\mathcal L^{X}(\pi)-\mathcal L^{Z}(\pi)\big|
\ \le\
\begin{cases}
2\,\beta\,L_\pi\,\Delta_R,&\text{case (1)},\\[2pt]
2\,L_R\,\Delta_R,&\text{case (2)},
\end{cases}
\qquad
\text{and}\qquad
\inf_\pi \mathcal L^{X}(\pi)-\inf_\pi \mathcal L^{Z}(\pi)\ \le\ \text{same bound}.
\]
\end{theorem}

\begin{proof}
Let $\varphi(t)=-\log\sigma(t)$, which is $1$–Lipschitz. Denote the DPO logit by
\[
\Lambda(X;a^w,a^\ell)\ :=\ \beta\!\left[\log\frac{\pi(a^w\mid X)}{\pi_{\mathrm{ref}}(a^w\mid X)}-\log\frac{\pi(a^\ell\mid X)}{\pi_{\mathrm{ref}}(a^\ell\mid X)}\right].
\]
Then
\[
\big|\mathcal L^{X}(\pi)-\mathcal L^{Z}(\pi)\big|
\ \le\ \E\big|\varphi(\Lambda(X;a^w,a^\ell))-\varphi(\Lambda(Z^{(R)}(X);a^w,a^\ell))\big|
\ \le\ \E\big|\Lambda(X;a^w,a^\ell)-\Lambda(Z^{(R)}(X);a^w,a^\ell)\big|.
\]
In case (1), oracle‑measurability lets us replace $X$ by $y=\fstar(X)$ and $Z^{(R)}(X)$ by $y'=\fstar(Z^{(R)}(X))$, and the Lipschitz condition yields
\[
\big|\Lambda(X;a^w,a^\ell)-\Lambda(Z^{(R)}(X);a^w,a^\ell)\big|\ \le\ \beta\,L_\pi\big(\metric(y,y')+\metric(y,y')\big)\ =\ 2\,\beta\,L_\pi\,\metric(y,y').
\]
Taking expectations gives the stated $2\,\beta\,L_\pi\,\Delta_R$ bound. In case (2), write the logit difference as $\beta^{-1}\!\big[(R_y(a^w)-R_{y'}(a^w))-(R_y(a^\ell)-R_{y'}(a^\ell))\big]$ and apply the reward‑Lipschitz bound twice to get $2\,L_R\,\metric(y,y')$; take expectations to conclude.
\end{proof}

\paragraph{Remarks.} (i) Because $\metric\ge0$, any display of the form $\E[\metric(\cdot,\cdot)]=0$ immediately implies $\metric(\cdot,\cdot)=0$ almost surely with respect to the summarizer's randomness, conditional on the realized subtree. We use the stronger almost‑sure formulation in the exact‐equivalence statements for clarity. (ii) If $\pi_{\mathrm{ref}}$ is \emph{not} oracle‑measurable, exact preservation of $\fstar$ alone does not guarantee equivalence of minimizers; the KL term can transmit presentational artifacts of $X$ into the optimizer. This is why the oracle‑indexed pair generator and oracle‑measurable $\pi_{\mathrm{ref}}$ hypothesis appears in Theorem~\ref{thm:dpo-exact}. (iii)~\emph{Formal verification status:} The companion Lean formalization in \texttt{FormalProofs/DPO.lean} provides machine-verified proofs for:
\begin{itemize}
  \item \textbf{Exact equivalence} (\texttt{dpo\_equivalence}, \texttt{dpo\_exact\_metric}): When local laws L1, L2, L3 hold exactly, DPO gap = 0. This is the primary theoretical result.
  \item \textbf{Crude bound} (\texttt{dpo\_gap\_oracle\_indexed}): Gap $\le 2\cdot M_{\mathrm{loss}}$ where $M_{\mathrm{loss}}$ bounds the loss function.
\end{itemize}
The specific Lipschitz bounds $2\beta L_\pi\Delta_R$ and $2L_R\Delta_R$ in Theorem~\ref{thm:dpo-gap} are \emph{not} machine-verified and should be treated as proof sketches. Formalizing these bounds with explicit Lipschitz constants is future work.

\section{Global Compatibility Variant and Algebraic Structure}\label{sec:global}

This section records a \emph{global} counterpart to the local, auditable laws used in the main text. The point is conceptual: if the two-route identities were to hold \emph{for all} pairs in $\Strings\times\Strings$, and if on-range parent calls depended only on child oracle values, then concatenation induces a \emph{merge law on oracle values}. All equalities below are stated in the task metric $\metric$: when we write that two oracle outputs are ``equal,'' this always means their $\metric$–distance is zero.

\paragraph{Formal verification.} The propositions in this section are machine-verified in \texttt{FormalProofs/GlobalAssumptions.lean}. The Lean formalization uses the notation A1, A2, A3 for the assumptions below, and provides the theorems \texttt{concat\_congruence}, \texttt{merge\_assoc}, and \texttt{merge\_identity}.

\begin{remark}[How to read ``equal at the oracle'']
We will say $y$ and $y'$ in $\Y$ are \emph{equal at the oracle} if $\metric(y,y')=0$. This can be thought of as identifying outcomes that the task metric cannot distinguish. More formally, this is the same as passing to the equivalence relation $y\sim y'$ iff $\metric(y,y')=0$ and reading all statements modulo $\sim$.
\end{remark}

\paragraph{Global laws.} We fix a \emph{deterministic} summarizer $g:\Strings\to\Strings$ and assume three global properties.

\begin{assumption}[Global $\fstar$–sufficiency]\label{ass:app1}
For every $z\in\Strings$,
\[
\metric\!\big(\fstar(g(z)),\,\fstar(z)\big)=0.
\]
\end{assumption}

\begin{assumption}[Global two-route compatibility]\label{ass:app2}
For every $u,v\in\Strings$,
\[
\metric\!\Big(\fstar(u\concat v),\ \fstar\!\big(g\big(g(u)\concat g(v)\big)\big)\Big)=0.
\]
\end{assumption}

\begin{assumption}[On-range merge depends only on child oracles]\label{ass:app3}
There exists a measurable map $M:\Y\times\Y\to\Y$ such that for all $u,v\in\Strings$,
\[
\metric\!\Big(\fstar\!\big(g\big(g(u)\concat g(v)\big)\big),\ M\big(\fstar(g(u)),\,\fstar(g(v))\big)\Big)=0,
\]
and $M$ is \emph{well-defined at the oracle level}: whenever $\metric(y_i,y_i')=0$ for $i=1,2$, then
\[
\metric\!\big(M(y_1,y_2),\,M(y_1',y_2')\big)=0.
\]
\end{assumption}

The first law says a one-pass summary preserves the task value (globally, not just on realized leaves). The second is the global version of the edgewise two-route equality, with the \emph{parent} summary explicit. The third law formalizes that, once inputs are on the range of $g$, the parent’s oracle depends only on the two child oracle values; it also says this dependence respects ``equality at the oracle'' (outputs cannot change when inputs are changed within $\metric$–zero).

\medskip

We now state two consequences. 


\begin{proposition}[Concatenation is a congruence for oracle equality on strings]\label{prop:congruence-strings}
\emph{(Lean: \texttt{concat\_congruence})}

If $a,a',b,b'\in\Strings$ satisfy $\metric\!\big(\fstar(a),\fstar(a')\big)=0$ and $\metric\!\big(\fstar(b),\fstar(b')\big)=0$, then
\[
\metric\!\big(\fstar(a\concat b),\ \fstar(a'\concat b')\big)=0.
\]
\end{proposition}

\begin{proof}
Let $U:=g\big(g(a)\concat g(b)\big)$ and $U':=g\big(g(a')\concat g(b')\big)$. By two applications of the triangle inequality,
\[
\metric\!\big(\fstar(a\concat b),\fstar(a'\concat b')\big)\ \le\ 
\underbrace{\metric\!\big(\fstar(a\concat b),\fstar(U)\big)}_{\text{(I)}}\ +\
\underbrace{\metric\!\big(\fstar(U),\fstar(U')\big)}_{\text{(II)}}\ +\
\underbrace{\metric\!\big(\fstar(U'),\fstar(a'\concat b')\big)}_{\text{(III)}}.
\]
By Assumption~\ref{ass:app2}, terms (I) and (III) are zero. For (II), Assumption~\ref{ass:app3} gives
\[
\metric\!\Big(\fstar(U),\,M\big(\fstar(g(a)),\fstar(g(b))\big)\Big)=0,\qquad
\metric\!\Big(\fstar(U'),\,M\big(\fstar(g(a')),\,\fstar(g(b'))\big)\Big)=0.
\]
Assumption~\ref{ass:app1} and the premises yield
\[
\metric\!\big(\fstar(g(a)),\fstar(g(a'))\big)\le \metric\!\big(\fstar(g(a)),\fstar(a)\big)+\metric\!\big(\fstar(a),\fstar(a')\big)+\metric\!\big(\fstar(a'),\fstar(g(a'))\big)=0,
\]
and analogously $\metric\!\big(\fstar(g(b)),\fstar(g(b'))\big)=0$. The oracle-level well-definedness in Assumption~\ref{ass:app3} therefore implies
\[
\metric\!\Big(M\big(\fstar(g(a)),\fstar(g(b))\big),\ M\big(\fstar(g(a')),\,\fstar(g(b'))\big)\Big)=0,
\]
which forces (II)$=0$ by triangle inequality. Hence the whole right-hand side is zero.
\end{proof}

\begin{proposition}[A merge law on oracle values (associative up to $\metric$)]\label{prop:odot}
\emph{(Lean: \texttt{merge\_assoc}, \texttt{merge\_identity})}

Define, for oracle values that actually arise on-range, the binary operation
\[
y_1\ \odot\ y_2\ :=\ M(y_1,y_2).
\]
Then for all $y_1,y_2,y_3$ in the on-range set $\fstar(g(\Strings))$, we have
\[
\metric\!\Big((y_1\odot y_2)\odot y_3,\ y_1\odot(y_2\odot y_3)\Big)=0,
\]
and $y_0:=\fstar(g(\emptystr))$ is a two-sided identity in the same $\metric$–sense:
\[
\metric(y_0\odot y,\ y)=\metric(y\odot y_0,\ y)=0\quad\text{for all on-range }y.
\]
Consequently, concatenation on $\Strings$ descends, via $\fstar$, to an associative merge on oracle values \emph{up to the task metric}.
\end{proposition}

\begin{proof}
Pick $a,b,c\in\Strings$ with $y_1=\fstar(g(a))$, $y_2=\fstar(g(b))$, $y_3=\fstar(g(c))$. By Assumption~\ref{ass:app3},
\[
(y_1\odot y_2)\odot y_3\ \text{ is oracle-equal to }\ \fstar\!\big(g\big(g\big(g(a)\concat g(b)\big)\ \concat\ g(c)\big)\big),
\]
and
\[
y_1\odot(y_2\odot y_3)\ \text{ is oracle-equal to }\ \fstar\!\big(g\big(g\big(a\big)\ \concat\ g\big(g(b)\concat g(c)\big)\big)\big).
\]
Applying Assumption~\ref{ass:app2} twice yields
\[
\metric\!\Big((y_1\odot y_2)\odot y_3,\ \fstar\big((a\concat b)\concat c\big)\Big)=0,\qquad
\metric\!\Big(y_1\odot (y_2\odot y_3),\ \fstar\big(a\concat(b\concat c)\big)\Big)=0.
\]
Since concatenation on $\Strings$ is literally associative, $(a\concat b)\concat c=a\concat(b\concat c)$, hence the two targets of the displays are equal and therefore at $\metric$–distance zero from one another. The identity statements follow in the same way with $b=c=\emptystr$ and Assumptions~\ref{ass:app1}–\ref{ass:app3}.
\end{proof}

\begin{remark}[Dropping the parent $g$ at a merge]
By Assumption~\ref{ass:app1}, $\metric\!\big(\fstar\!\big(g(g(u)\concat g(v))\big),\,\fstar\!\big(g(u)\concat g(v)\big)\big)=0$ for all $u,v$. Thus, in Assumption~\ref{ass:app2}, one may replace the right-hand side by $\fstar\!\big(g(u)\concat g(v)\big)$ without changing any conclusion in the $\metric$–sense. We keep the explicit parent call to mirror the deployed pipeline.
\end{remark}

\paragraph{What this variant does \emph{not} assume.} We have not assumed that $M$ is literally associative on \emph{all} of $\Y$, nor that $\fstar$ embeds $\Strings$ into $\Y$ as a strict monoid. Propositions~\ref{prop:congruence-strings}–\ref{prop:odot} assert the exact strength warranted by the applications: along on-range values produced by $g$, and with equalities read in the task metric, the two-route global laws make concatenation behave as an associative merge on oracle values. This global structure can be useful when the task semantics truly admit a merge law (e.g., counts); the main-text guarantees, however, require only the local, auditable laws.




\section{Information-Theoretic View of OPS}\label{sec:info-sieve}

This appendix makes explicit the connection between Oracle-Preserving Summarization (OPS) and the \emph{Information Sieve} of \citet{steeg2016informationsieve}. The sieve decomposes an input $X$ into a latent factor $Y$ and \emph{remainder information} $\bar X$ such that $(Y,\bar X)$ reconstructs $X$ exactly while $Y$ captures as much multivariate dependence (Total Correlation) as possible and $\bar X$ is independent of $Y$. OPS instantiates the same architecture, but the objective is supervised: rather than maximizing dependence within $X$, we require that the retained summary $Z=g(X)$ be a sufficient statistic for the oracle $f^*(X)$.

\begin{theorem}[Oracle information decomposition]\label{thm:info-decomp}
Let $X$ be a block or node in the hierarchy, let $Z=g(X)$ be its summary, and let $\nu$ denote any (possibly stochastic) remainder such that $X$ is measurable with respect to $(Z,\nu)$. Then
\[
I(X; f^*) \;=\; I(Z; f^*) \;+\; I(\nu; f^* \mid Z).
\]
In particular, $Z$ preserves all oracle information if and only if $I(\nu; f^* \mid Z)=0$.
\end{theorem}

\begin{proof}
Apply the chain rule for mutual information with $X = (Z,\nu)$.
\end{proof}

Condition~(C1) (Sufficiency) enforces $I(\nu; f^* \mid Z)=0$ at every realized leaf by checking that the oracle agrees on $b$ and $g(b)$. Condition~(C2) (Merge Consistency) ensures that when $Z$ is formed by hierarchically merging child summaries $(Z_{u_L},Z_{u_R})$, the remainder created at that merge is also orthogonal to $f^*$. Combining Theorem~\ref{thm:info-decomp} with the incremental construction of the Information Sieve (Theorem~2.1 of \citet{steeg2016informationsieve}) yields an additive view of the hierarchy: each merge acts as a sieve layer whose discarded bits are independent of the oracle. Consequently, if the local audits certify that $I(\nu; f^* \mid Z)=0$ on every edge, the root summary preserves the full mutual information $I(X_{\text{root}}; f^*)$ present in the original document.

Finally, Section~5 of \citet{steeg2016informationsieve} observes that discarding the remainder induces a lossy compression of the raw input. OPS inherits this interpretation: summaries $g(x)$ need not suffice to reconstruct the surface text, but they are \emph{lossless with respect to $f^*$}. The ``loss'' lives entirely in the nuisance remainder, which by design is independent of the oracle and therefore irrelevant for downstream coding tasks.

\section{Relationship to Mergeable Summaries}\label{app:mergeable}

This appendix formalizes the connection between our Oracle-Preserving Summarization (OPS) framework and the theory of \emph{mergeable summaries} developed for streaming and distributed databases \citep{Agarwal2013MergeableTODS}. We show that OPS strictly generalizes mergeable summaries: when the oracle $\fstar$ is symmetric and the summarizer $g$ is deterministic, our Parentwise Compatibility law~\eqref{eq:leaf_compatibility} is identical to the mergeability condition. When $\fstar$ is order-sensitive or $g$ is stochastic, OPS extends the notion of mergeability to the free monoid of strings by replacing hand-designed merge operators with learned semantic merges implemented by an LLM.

\subsection{Formal Setup of Mergeable Summaries}

In the classical setting, data arrives as a collection of items, and the objective is to maintain a compact synopsis that can be merged efficiently. Let $\mathcal{D}$ denote datasets (typically multisets) and $\mathcal{S}$ the sketch space with $|\mathcal{S}|\ll|\mathcal{D}|$.
\begin{itemize}
    \item \textbf{Data:} $D_1, D_2 \in \mathcal{D}$.
    \item \textbf{Operation:} Multiset union $\uplus$, which is associative and commutative.
    \item \textbf{Summary:} A function $S:\mathcal{D}\to\mathcal{S}$.
    \item \textbf{Query:} $Q$ maps a sketch to the target statistic, with distortion bounded by $\epsilon$.
\end{itemize}

\begin{definition}[Mergeable Summary]
A summary $S$ is mergeable if there exists $\mathcal{M}:\mathcal{S}\times\mathcal{S}\to\mathcal{S}$ such that for all $D_1,D_2$,
\[
    d_{\Y}\!\left(Q\!\left(\mathcal{M}(S(D_1),S(D_2))\right),\, Q(D_1\uplus D_2)\right) \le \epsilon.
\]
\end{definition}

The guarantee holds uniformly over arbitrary merge trees, allowing sketches to be composed without revisiting the raw inputs.

\subsection{Mapping OPS to Mergeable Summaries}

Table~\ref{tab:mergeable_mapping} aligns the notation used in Section~\ref{sec:mechanism} with the streaming literature. The distinction is that OPS works over the free monoid of strings $(\Strings,\concat)$, where $\concat$ is generally non-commutative, and uses an LLM summarizer $g$ instead of a fixed algebraic operator.

\begin{table}[t]
    \centering
    \begin{tabular}{p{0.23\linewidth}p{0.33\linewidth}p{0.33\linewidth}}
        \toprule
        \textbf{Component} & \textbf{OPS (this paper)} & \textbf{Mergeable summaries} \\
        \midrule
        Input & Strings $u,v\in\Strings$ & Datasets $D_1,D_2\in\mathcal{D}$ \\
        Combination & Concatenation $\concat$ (non-commutative) & Union $\uplus$ (commutative) \\
        Summarizer & $g:\Strings\rightsquigarrow\Strings$ & $S:\mathcal{D}\to\mathcal{S}$ \\
        Task/query & Oracle $\fstar:\Strings\to\Y$ & Query $Q$ \\
        Merge algorithm & $g(g(u)\concat g(v))$ & $\mathcal{M}(S(D_1),S(D_2))$ \\
        Consistency law & Parentwise compatibility \eqref{eq:leaf_compatibility} & Mergeability condition \\
        \bottomrule
    \end{tabular}
    \caption{Mapping OPS to the mergeable summaries framework.}
    \label{tab:mergeable_mapping}
\end{table}

\subsection{Equivalence in the Commutative Case}

Suppose the task ignores order, such as counting keyword occurrences or extracting the set of unique actors. In this case, $\fstar$ acts as a homomorphism from $(\Strings,\concat)$ to a commutative monoid $(\Y,+)$:
\[
    \fstar(u\concat v) = \fstar(u) + \fstar(v),
\]
up to the oracle metric $\metric$.

Parentwise compatibility then states that re-summarizing the child summaries is inert at the oracle:
\[
    \E\!\Big[\metric\!\Big(\fstar\big(g(g(u)\concat g(v))\big),\, \fstar(g(u)\concat g(v))\Big)\Big]=0.
\]
When $g$ is deterministic, the expectation drops and the condition becomes pointwise.

\begin{proposition}[Reduction to Mergeable Sketches]
\label{prop:mergeable_reduction}
Let $\fstar$ be commutative so that $\fstar(u\concat v)=\fstar(v\concat u)$, and let $g$ be deterministic. If $g$ satisfies exact Parentwise Compatibility~\eqref{eq:leaf_compatibility} on every realized merge, then $g$ is a mergeable summary for $\fstar$ in the sense of \citet{Agarwal2013MergeableTODS}.
\end{proposition}

\begin{proof}
Define the merge algorithm $\mathcal{M}(s_1,s_2)=g(s_1\concat s_2)$, where $s_i=g(u_i)$ are summaries. For any $u,v$ we compare two routes: (i) concatenate the summaries $s_u,s_v$ and apply $\fstar$, yielding $\fstar(g(u)\concat g(v))$; (ii) merge the summaries via $\mathcal{M}$ and apply $\fstar$, yielding $\fstar(g(g(u)\concat g(v)))$. Parentwise compatibility enforces equality between these two values, so the mergeability condition holds with $\epsilon=0$. When the oracle is commutative and Condition~\eqref{eq:leaf_sufficiency} holds on the leaves, the concatenated summaries themselves agree with the concatenated raw text under $\fstar$, recovering the classical guarantee.
\end{proof}

Note that, unlike classical sketches where data arrive already as sketch inputs, Condition~(C2) in our framework also constrains the \emph{raw} concatenations $u\concat v$ to agree with their summarized counterparts. The parsing step from text to summary is therefore part of the algebra itself, which makes the OPS guarantees strictly stronger than standard mergeability claims.

\begin{remark}[The universal merge]\label{rem:universal_merge}
Classical sketches design bespoke $\mathcal{M}$ operators (vector addition for Count-Min, rank pruning for quantile sketches). OPS instead uses a universal merge: concatenate the summaries and apply the same summarizer again. The summarizer learns the algebra needed by the downstream oracle during its initial reduction pass, so no task-specific merge code is required.
\end{remark}

\subsection{Generalization to Non-Commutative Tasks}

The OPS framework extends mergeability to the free monoid where $u\concat v\neq v\concat u$. Narrative queries---``Did the protest in paragraph $A$ trigger the crackdown in paragraph $B$?''---depend on order and context. Standard sketches over multisets cannot express these dependencies because the union operator discards order. Parentwise compatibility enforces a non-commutative analogue: the interaction between summarized blocks mimics the interaction between the original texts when evaluated by $\fstar$. This is precisely the condition audited in Figure~\ref{fig:local_compatibility}.

\subsection{Error Propagation and Bounds}

Mergeable summaries typically admit error bounds that grow logarithmically with tree height. Theorems~\ref{thm:one-pass} and~\ref{thm:multi-round} give an analogous ``zero-error'' propagation: if local merges satisfy Parentwise Compatibility exactly (in expectation), the global error collapses to zero regardless of tree shape. When local failures are bounded by $\epsilon$, our deployment bound~\eqref{eq:error_budget} plays the role of classical error accumulation: the total distortion is controlled by a weighted sum of leaf, merge, and stabilization failure rates. This is the semantic counterpart to bounding sketch errors with respect to the number of merge operations; it justifies the probabilistic audit in Section~\ref{sec:audit}, where random sampling over the summary tree produces a high-confidence certificate for the entire hierarchy.

\section{Existence of Oracle-Preserving Summaries}
\label{app:existence}

Two existence notions matter: set-theoretic (\emph{some} $g$ exists) and operational (a \emph{short} $g$ exists that fits a budget). We state both in terms of the task metric.

\begin{proposition}[Canonical representative map in metric form]
\label{prop:canonical}
Because $\Strings$ is countable, define $x\sim x'$ iff $\metric(\fstar(x),\fstar(x'))=0$. Fix a total order on $\Strings$ and let $c:\Y\to\Strings$ send $y$ to the least $x$ with $\metric(\fstar(x),y)=0$. Then
\[
g_{\mathrm{can}}(x):=c\big(\fstar(x)\big)
\quad\text{satisfies}\quad
\metric\!\big(\fstar(g_{\mathrm{can}}(x)),\,\fstar(x)\big)=0\ \ \text{for all }x,
\]
and $g_{\mathrm{can}}$ is idempotent on its range. Hence $g_{\mathrm{can}}\in\Gstar$ in the sense of Definition~\ref{def:exact-suff-stable}.
\end{proposition}

\begin{proposition}[Compact encodings]
\label{prop:encoding}
If there exists an injective encoder $\mathrm{enc}:\Y\to\Strings$ of uniformly bounded length such that
$\metric(\fstar(\mathrm{enc}(y)),y)=0$ for all $y\in\Y$, then
\[
g_{\mathrm{enc}}(x):=\mathrm{enc}(\fstar(x))
\quad\text{obeys}\quad
\metric\!\big(\fstar(g_{\mathrm{enc}}(x)),\,\fstar(x)\big)=0
\]
and is idempotent on its range. Thus $g_{\mathrm{enc}}\in\Gstar$ and returns short strings.
\end{proposition}

Proposition~\ref{prop:canonical} certifies internal consistency: exact, idempotent, oracle-preserving normal forms always exist; Proposition~\ref{prop:encoding} records the operational version used in practice (the summary \emph{is} a compact encoding of the oracle).




\section{Worked Examples and Templates}
\label{sec:examples}

This section provides instantiation templates for four common task types, showing how to check Conditions~\eqref{eq:leaf_sufficiency}--\eqref{eq:leaf_idempotence} in practice. Each example specifies the oracle space $\Y$, the metric $\metric$, and concrete audit procedures for leaf sufficiency, merge consistency, and on-range idempotence. The first three examples cover standard deployment scenarios (binary event detection, structured entity extraction, and numeric aggregation); the fourth is a minimal counterexample demonstrating that on-range stability (Condition~\eqref{eq:leaf_idempotence}) cannot be omitted without breaking multi-round preservation.

\subsection{Binary event detection (QA/passage-level)}
\begin{itemize}
  \item $Y=\{0,1\}$, $\metric(y,y')=\mathbf 1\{y\neq y'\}$.
  \item Sufficiency check: ask $g$ to emit \texttt{EVENT: YES/NO} for each block; verify $\E[\mathbf 1\{\fstar(g(b))\neq \fstar(b)\}]=0$ within CI.
  \item Merge-consistency check: for realized $u$, compare $\fstar(g(u_L)\concat g(u_R))$ to $\fstar(g(g(u_L)\concat g(u_R)))$; failures indicate cross‑block evidence interactions not preserved by $g$.
\end{itemize}

\subsection{Entity extraction (tuple)}
\begin{itemize}
  \item $Y\subseteq\mathcal A\times\mathcal{ACT}\times\mathcal T\times\mathcal P$ with weighted Hamming distance on fields.
  \item Canonicalization: require one actor/action per line in fixed order; idempotence checks use the same schema to keep $g$ on‑range.
\end{itemize}

\subsection{Counting (numeric aggregation)}
\begin{itemize}
  \item $Y=\mathbb N$, $\metric(y,y')=|y-y'|$.
  \item Global variant instantiates $y_1\odot y_2=y_1\concat y_2$ under exact $\fstar$‑sufficiency; edgewise audits reduce to additivity checks at realized merges.
\end{itemize}

\subsection{Stability Is Substantive: A Minimal Counterexample}
\label{app:stability}

Let $\Strings$ contain special tokens \textsc{pos}, \textsc{neg}, and let $\fstar(\textsc{pos})=1$, $\fstar(\textsc{neg})=0$, with general $\fstar(x)\in\{0,1\}$. Define
\[
g(x)=
\begin{cases}
\textsc{pos},& \fstar(x)=1,\,x\notin\{\textsc{pos},\textsc{neg}\},\\
\textsc{neg},& \fstar(x)=0,\,x\notin\{\textsc{pos},\textsc{neg}\},\\
\textsc{neg},& x\in\{\textsc{pos},\textsc{neg}\}.
\end{cases}
\]
Then \ref{eq:leaf_sufficiency} holds on leaves never equal to \{\textsc{pos},\textsc{neg}\}, but $g$ is not on‑range stable since $\fstar(g(\textsc{pos}))=0\neq 1=\fstar(\textsc{pos})$. Thus additional normalization rounds can flip the label; \ref{eq:leaf_idempotence} is required to guarantee Theorem~\ref{thm:multi-round}.

